\begin{trapbox}

\begin{description}[style=nextline,leftmargin=2.8em,labelsep=0.5em,itemsep=0.4em]
 \item[\textbf{\textcolor{CTrap}{易错点一：把双曲线公式写成不带绝对值的椭圆式}}]
 错写 $|PF_1|=a+ex_0$、$|PF_2|=a-ex_0$。
 \textbf{纠正：} 双曲线两支都存在，$e>1$ 且 $|x_0|\ge a$：
 右支 $a-ex_0<0$，左支 $a+ex_0<0$。
 必须统一写 $|PF_1|=|a+ex_0|$、$|PF_2|=|a-ex_0|$。

 \item[\textbf{\textcolor{CTrap}{易错点二：左右焦点与同侧焦点混淆}}]
 $F_1$ 是\textbf{左焦点}，并不永远代表“远焦点”。
 \textbf{纠正：} 右支 $F_1$ 远、$F_2$ 近；左支 $F_1$ 近、$F_2$ 远。
 直接用“\textbf{同侧短 $=e|x_0|-a$，异侧长 $=e|x_0|+a$}”辨别。

 \item[\textbf{\textcolor{CTrap}{易错点三：只看横坐标符号，却未确认点在双曲线上}}]
 写出 $x_0=2$ 就认为一定在右支是不对的（例如 $a=4$ 时）。
 \textbf{纠正：} 曲线上的点必须满足 $|x_0|\ge a$，且满足原方程；
 $-a<x_0<a$ 内根本没有双曲线点。

 \item[\textbf{\textcolor{CTrap}{易错点四：把 $c^2$ 写成 $a^2-b^2$}}]
 \textbf{纠正：} 双曲线是 $c^2=a^2+b^2$，
 因而 $e=c/a>1$；椭圆才是 $c^2=a^2-b^2$。

 \item[\textbf{\textcolor{CTrap}{易错点五：平方开方后丢掉绝对值}}]
 $\sqrt{(ex_0-a)^2}=|ex_0-a|$，不能直接写成 $ex_0-a$。
 \textbf{纠正：} 先保留绝对值，或先根据所在支证明括号非负，再去掉它。

 \item[\textbf{\textcolor{CTrap}{易错点六：只凭 $|x_0|\ge a$ 就认为任意 $y_0$ 都可用}}]
 例如 $x_0=5$、$a=4$ 时，$y_0$ 仍须满足曲线方程。
 \textbf{纠正：} $|x_0|\ge a$ 是\textbf{必要条件}，不是任意给定点在曲线上的充分条件。

 \item[\textbf{\textcolor{CTrap}{易错点七：记混椭圆与双曲线的一变量设法}}]
 椭圆两段\textbf{和}为 $2a$，因此设 $a+t,a-t$；
 双曲线两段\textbf{差的绝对值}为 $2a$，因此设 $t+a,t-a$。
 \textbf{纠正：} 先看是“和固定”还是“差固定”，再决定怎样设。

 \item[\textbf{\textcolor{CTrap}{易错点八：换成纵轴焦点时仍用 $x_0$}}]
 对 $\frac{y^2}{a^2}-\frac{x^2}{b^2}=1$，
 若 $F_1(0,-c),F_2(0,c)$，应改成
 $|PF_1|=|a+ey_0|$、$|PF_2|=|a-ey_0|$。
\end{description}

\medskip
\textbf{\textcolor{CTrap}{考场 10 秒自检}}
\[
 \boxed{\bigl||PF_1|-|PF_2|\bigr|=2a},\qquad
 \boxed{|PF_1|+|PF_2|=2e|x_0|}.
\]
若任何一式不成立，先检查焦点左右、横坐标符号与绝对值。

\end{trapbox}
